\chapter{Integrais de linha e integrais múltiplas}\label{ch:b2:multint}

Este capítulo estende a integração de intervalos a curvas e a
domínios do plano e do espaço. As \omterm{def:b2:multint:lineint}{integrais de linha} integram uma
\emph{\omterm{def:b2:multint:lineint}{forma diferencial}} $P\,\dd x + Q\,\dd y$ ao longo de um
arco orientado; as integrais duplas e triplas integram funções sobre
domínios de dimensão dois e três. As duas teorias se encontram no
\emph{teorema de Green--Riemann}, o teorema fundamental do cálculo
em dimensão dois, e a principal ferramenta de cálculo em todo o
capítulo é a \emph{fórmula de mudança de variáveis}, cujo fator de distorção é
o valor absoluto do \omterm{def:b2:linalg:det}{determinante} jacobiano.

\section{Integrais de linha}

\begin{definition}[Forma diferencial; integral de linha]\label{def:b2:multint:lineint}
Seja $U \subseteq \R^2$ um \omterm{def:b2:metric:topology}{aberto}. Uma \emph{forma
diferencial}\index{forma diferencial} de
grau $1$ e classe $\mathcal{C}^0$ em $U$ é uma expressão
$\omega = P\,\dd x + Q\,\dd y$ com $P, Q \colon U \to \R$
\omterm{def:b2:metric:continuity}{contínuas} --- formalmente, uma aplicação \omterm{def:b2:metric:continuity}{contínua} de $U$ no \omterm{def:b2:linalg:dual}{dual}
de $\R^2$, $\omega(M) = P(M)\,e_1^* + Q(M)\,e_2^*$. Para um
arco $\gamma \colon [a, b] \to U$ de classe $\mathcal{C}^1$, $\gamma(t) =
(x(t), y(t))$, a \emph{integral de linha}\index{integral de linha} de
$\omega$ ao longo de $\gamma$
é
\[
\int_\gamma \omega
= \int_a^b \Bigl(P(\gamma(t))\,x'(t)
+ Q(\gamma(t))\,y'(t)\Bigr)\,\dd t .
\]
As definições se estendem literalmente a $\R^3$ (formas $P\,\dd x +
Q\,\dd y + R\,\dd z$) e a arcos $\mathcal{C}^1$ por pedaços (soma
sobre os pedaços).
\end{definition}

\begin{proposition}[Invariância e orientação]\label{prop:b2:multint:lineinv}
A \omterm{def:b2:multint:lineint}{integral de linha} não muda sob uma \omterm{def:b2:curves:reparam}{mudança de parâmetro}
$\mathcal{C}^1$ crescente, e muda de sinal sob uma decrescente. Ela
depende portanto apenas do \omterm{def:b2:curves:reparam}{arco geométrico} \emph{orientado}.
\end{proposition}

\begin{proof}
Se $\theta \colon [c, d] \to [a, b]$ é uma \omterm{def:b2:curves:reparam}{mudança de parâmetro} e
$\tilde\gamma = \gamma \circ \theta$, então, pela regra da cadeia e pela
mudança de variáveis a uma variável $t = \theta(u)$,
\[
\int_{\tilde\gamma}\omega
= \int_c^d \bigl(P(\gamma(\theta(u)))\,x'(\theta(u))
+ Q(\gamma(\theta(u)))\,y'(\theta(u))\bigr)\,\theta'(u)\,\dd u
= \pm\int_a^b \bigl(Px' + Qy'\bigr)(t)\,\dd t ,
\]
com sinal $+$ se $\theta$ é crescente ($\theta(c) = a$) e $-$
se for decrescente (os limites se trocam).
\end{proof}

\begin{example}[Trabalho de uma força; circulação]\label{ex:b2:multint:work}
Se $F = (P, Q)$ é um campo de forças, $\int_\gamma P\dd x + Q\dd y =
\int_a^b \langle F(\gamma(t)), \gamma'(t)\rangle\,\dd t$ é o
\emph{trabalho} de $F$ ao longo de $\gamma$. Para $\omega = -y\,\dd x +
x\,\dd y$ ao longo do círculo unitário $\gamma(t) =
(\cos t, \sin t)$ percorrido no sentido anti-horário:
\[
\int_\gamma \omega
= \int_0^{2\pi}\bigl((-\sin t)(-\sin t)
+ \cos t\cos t\bigr)\,\dd t = 2\pi ,
\]
o dobro da \omterm{def:b2:multint:domain}{área} encerrada --- um primeiro indício de Green--Riemann.
\end{example}

\begin{example}[Uma integral, duas parametrizações, uma
armadilha de sinal]\label{ex:b2:multint:semicircle}
Calcule $\int_\gamma x\,\dd y$ ao longo do semicírculo unitário
superior de $(1, 0)$ a $(-1, 0)$. Com $\gamma(t) =
(\cos t, \sin t)$, $t \in \intcc0\pi$:
\[
\int_0^\pi\cos t\cdot\cos t\,\dd t = \frac\pi2 .
\]
Com a parametrização por gráfico $x \mapsto (x, \sqrt{1 -
x^2})$, $x$ indo de $1$ a $-1$ (atenção ao sentido!):
\[
\int_1^{-1}x\cdot\frac{-x}{\sqrt{1 - x^2}}\,\dd x
= \int_{-1}^{1}\frac{x^2}{\sqrt{1 - x^2}}\,\dd x
= \frac\pi2
\]
($x = \sin u$ a reduz a uma \omterm{pb:b2:multint:1}{integral de Wallis}). Mesmo valor,
como a~\cref{prop:b2:multint:lineinv} garante --- mas só
porque os dois percursos vão de $(1,0)$ a $(-1,0)$; inverter o
sentido troca o sinal. Fechar o caminho ao longo do eixo $x$
(em que $\dd y = 0$) não acrescenta nada, e o total $\frac\pi2$
é a \omterm{def:b2:multint:domain}{área} do semidisco: a primeira instância das fórmulas de
\omterm{def:b2:multint:domain}{área} pelo bordo de Green--Riemann abaixo.
\end{example}

\begin{definition}[Formas exatas e fechadas]\label{def:b2:multint:exact}
A forma $\omega = P\,\dd x + Q\,\dd y$ de classe $\mathcal{C}^0$
é \emph{exata}\index{forma exata} em $U$ se existe $f \in \mathcal{C}^1(U)$
(um \emph{potencial}\index{potencial}) com $\omega = \dd f$, isto é, $P = f_x$ e $Q
= f_y$. Uma forma $\mathcal{C}^1$ é \emph{fechada}\index{forma fechada} se $P_y = Q_x$ em
$U$.
\end{definition}

\begin{theorem}[Teorema fundamental para integrais de
linha]\label{thm:b2:multint:ftc}
Se $\omega = \dd f$ é \omterm{def:b2:multint:exact}{exata} e $\gamma$ é um arco $\mathcal{C}^1$
por pedaços em $U$ de $A$ a $B$, então
\[
\int_\gamma \omega = f(B) - f(A) .
\]
Em particular, a integral de uma \omterm{def:b2:multint:exact}{forma exata} ao longo de qualquer arco fechado
é nula, e toda forma $\mathcal{C}^1$ \omterm{def:b2:multint:exact}{exata} é fechada.
\end{theorem}

\begin{proof}
$\frac{\dd}{\dd t}f(\gamma(t)) = f_x(\gamma(t))x'(t) +
f_y(\gamma(t))y'(t)$ pela regra da cadeia
(\cref{ch:b2:diffcalc}), de modo que o integrando em
\cref{def:b2:multint:lineint} é a derivada de $t \mapsto
f(\gamma(t))$, e o teorema fundamental do cálculo dá o
resultado em cada pedaço; os valores intermediários se telescopam. Que as
formas $\mathcal{C}^1$ \omterm{def:b2:multint:exact}{exatas} sejam fechadas é o teorema de Schwarz: $P_y =
f_{xy} = f_{yx} = Q_x$.
\end{proof}

\begin{example}[Reconstruindo um potencial]\label{ex:b2:multint:potential}
Seja $\omega = y\,\eu^{xy}\,\dd x + (x\,\eu^{xy} + 2y)\,\dd y$
em $\R^2$. Ela é fechada: ambas as derivadas cruzadas valem
$\eu^{xy}(1 + xy)$. Para achar um \omterm{def:b2:multint:exact}{potencial}, integre $P$ em
$x$ com $y$ fixo:
\[
f(x, y) = \int y\,\eu^{xy}\,\dd x = \eu^{xy} + c(y),
\]
depois ajuste $c$ igualando $f_y$: $x\,\eu^{xy} + c'(y) =
x\,\eu^{xy} + 2y$ dá $c(y) = y^2$. Assim $f(x,y) = \eu^{xy} +
y^2$ e, para qualquer arco $\mathcal C^1$ por pedaços de $(0,0)$
a $(1,1)$,
\[
\int_\gamma\omega = f(1,1) - f(0,0) = (\eu + 1) - 1 = \eu ,
\]
independentemente do caminho --- a receita em dois passos (integrar
em $x$, corrigir em $y$) é a recíproca prática do
\cref{thm:b2:multint:ftc} em domínios em que as \omterm{def:b2:multint:exact}{formas fechadas} são
\omterm{def:b2:multint:exact}{exatas}.
\end{example}

\begin{example}[Fechada não implica exata]\label{ex:b2:multint:angleform}
Em $U = \R^2 \setminus \{0\}$, a \emph{forma ângulo}
\[
\omega = \frac{-y\,\dd x + x\,\dd y}{x^2 + y^2}
\]
é fechada (cálculo direto: tanto $P_y$ quanto $Q_x$ valem
$\frac{y^2 - x^2}{(x^2+y^2)^2}$), mas sua integral ao longo do círculo
unitário vale $2\pi \neq 0$ (mesmo cálculo do
\cref{ex:b2:multint:work}, dividido por $1$): $\omega$ não é \omterm{def:b2:multint:exact}{exata}
em $U$. Localmente, $\omega = \dd\theta$ para uma determinação
$\theta$ do ângulo polar; a falha é global --- o ângulo
não pode ser definido \omterm{def:b2:metric:continuity}{continuamente} em torno do furo. Em domínios
sem buracos a patologia desaparece: num \omterm{def:b2:metric:topology}{aberto} \emph{estrelado},
toda forma $\mathcal{C}^1$ fechada é \omterm{def:b2:multint:exact}{exata}
(lema de Poincaré, \cref{exo:b2:multint:8}).
\end{example}

\section{Integrais duplas}

Tomamos por adquirida a teoria da integral de Riemann a uma
variável (volume do primeiro ano de graduação e \cref{ch:b2:integration}) e esboçamos
sua versão a duas variáveis. Uma função $f$ \omterm{def:b2:metric:continuity}{contínua} num retângulo
$R = [a, b] \times [c, d]$ tem uma integral dupla $\iint_R f$,
definida por somas de Riemann sobre grades exatamente como a uma variável, e
calculada por iteração:

\begin{theorem}[Fubini num retângulo]\label{thm:b2:multint:fubini}
Para $f$ \omterm{def:b2:metric:continuity}{contínua} em $R = [a,b] \times [c,d]$,
\[
\iint_R f
= \int_a^b \Bigl(\int_c^d f(x, y)\,\dd y\Bigr)\dd x
= \int_c^d \Bigl(\int_a^b f(x, y)\,\dd x\Bigr)\dd y .
\]
\end{theorem}

\begin{proof}
Ponha $F(x) = \int_c^d f(x, y)\,\dd y$. A \omterm{def:b2:metric:continuity}{continuidade} uniforme de $f$ no
\omterm{def:b2:metric:compact}{compacto} $R$ torna $F$ \omterm{def:b2:metric:continuity}{contínua} (estimativa dominada:
$\abs{F(x) - F(x')} \leq (d - c)\sup_y\abs{f(x,y) - f(x',y)}$).
Subdivida agora $[a,b]$ e $[c,d]$ em $n$ partes iguais, o que dá uma
grade de células $R_{ij}$ de \omterm{def:b2:multint:domain}{área} $\Delta x\,\Delta y$. Em cada célula,
$\inf_{R_{ij}} f \cdot \Delta x \Delta y \leq
\int_{x_{i-1}}^{x_i}\int_{y_{j-1}}^{y_j} f(x,y)\,\dd y\,\dd x \leq
\sup_{R_{ij}} f \cdot \Delta x \Delta y$ pela monotonia da
integral a uma variável (aplicada duas vezes). Somando sobre as células, a
integral iterada $\int_a^b F$ fica espremida entre as somas de Riemann
inferior e superior da grade; pela \omterm{def:b2:metric:continuity}{continuidade} uniforme, ambas as somas
convergem ao valor comum que define $\iint_R f$ quando $n \to
\infty$. O mesmo argumento se aplica com os papéis de $x$ e $y$
trocados, de modo que ambas as integrais iteradas valem $\iint_R f$.
\end{proof}

\begin{remark}
A \omterm{def:b2:metric:continuity}{continuidade} num retângulo \omterm{def:b2:metric:compact}{compacto} faz trabalho de verdade na
demonstração de Fubini: ela fornece a \omterm{def:b2:metric:continuity}{continuidade} uniforme que
espreme as somas de Riemann. Para integrandos mais selvagens o
enunciado realmente falha --- há funções cujas duas
integrais iteradas existem e diferem. O teorema geral honesto,
com a integrabilidade como única hipótese, é o
teorema de Fubini para a integral de Lebesgue, demonstrado no
volume do terceiro ano de graduação; tudo neste capítulo permanece no
quadro \omterm{def:b2:metric:continuity}{contínuo} em que a demonstração elementar acima é
\omterm{def:b2:metric:complete}{completa}.
\end{remark}

\begin{definition}[Domínios elementares]\label{def:b2:multint:domain}
Um domínio $D \subseteq \R^2$ é \emph{elementar em $y$} se
\[
D = \{(x, y) : a \leq x \leq b,\
\varphi_1(x) \leq y \leq \varphi_2(x)\}
\]
com $\varphi_1 \leq \varphi_2$ \omterm{def:b2:metric:continuity}{contínuas} em $[a,b]$
(elementar em $x$: simetricamente). Para $f$ \omterm{def:b2:metric:continuity}{contínua} num
$D$ elementar em $y$,
\[
\iint_D f
= \int_a^b\Bigl(
\int_{\varphi_1(x)}^{\varphi_2(x)} f(x,y)\,\dd y\Bigr)\dd x ,
\]
e verifica-se (estendendo $f$ por um argumento de aproximação, ou
subdividindo) que, quando $D$ é elementar nos dois sentidos, as
duas integrais iteradas coincidem. Domínios recortados em finitos
pedaços elementares são tratados por aditividade, e a
\emph{área}\index{área!de um domínio plano} de $D$ é $\operatorname
{Area}(D) = \iint_D 1$.
\end{definition}

\begin{example}\label{ex:b2:multint:triangle}
No triângulo $D = \{0 \leq x \leq 1,\ 0 \leq y \leq x\}$:
\[
\iint_D xy \,\dd x\,\dd y
= \int_0^1 x\Bigl(\int_0^x y\,\dd y\Bigr)\dd x
= \int_0^1 x\cdot\frac{x^2}{2}\,\dd x = \frac18 .
\]
Trocando a ordem ($x$ de $y$ a $1$): $\int_0^1
y\bigl(\int_y^1 x\,\dd x\bigr)\dd y = \int_0^1
y\,\frac{1 - y^2}{2}\,\dd y = \frac18$ --- mesmo valor, cálculo
diferente: escolher bem a ordem de integração é metade do
ofício.
\end{example}

\begin{example}[Quando só uma ordem funciona]\label{ex:b2:multint:swaporder}
Calcule $I = \displaystyle\int_0^1\!\!\int_x^1
\eu^{y^2}\,\dd y\,\dd x$. Como está escrita, a integral interna
$\int\eu^{y^2}\dd y$ não tem primitiva elementar: o
cálculo trava. Mas o domínio é o triângulo $0 \leq
x \leq y \leq 1$, que é elementar nos dois sentidos;
trocando a ordem,
\[
I = \int_0^1\!\!\int_0^y \eu^{y^2}\,\dd x\,\dd y
= \int_0^1 y\,\eu^{y^2}\,\dd y
= \Bigl[\tfrac12\eu^{y^2}\Bigr]_0^1 = \frac{\eu - 1}{2} .
\]
A variável interna $x$ não aparecia em lugar algum do integrando, de modo que
integrá-la primeiro produziu exatamente o fator $y$ que
torna imediata a integral externa. Moral: Fubini não é apenas
uma licença para iterar --- é uma licença para \emph{escolher},
e a ordem certa pode transformar uma integral impossível numa
única linha. Esboce sempre o domínio e leia as duas
descrições antes de começar.
\end{example}

\begin{theorem}[Mudança de variáveis]\label{thm:b2:multint:changevar}
Seja $\Phi \colon U' \to U$ um difeomorfismo $\mathcal{C}^1$
entre \omterm{def:b2:metric:topology}{abertos} de $\R^2$, seja $K \subseteq U$ um domínio
\omterm{def:b2:metric:compact}{compacto} recortado em pedaços elementares com $K' = \Phi^{-1}(K)$, e
seja $f$ \omterm{def:b2:metric:continuity}{contínua} em $K$. Então
\[
\iint_K f(x, y)\,\dd x\,\dd y
= \iint_{K'} f\bigl(\Phi(u, v)\bigr)\,
\abs{\det J_\Phi(u, v)}\,\dd u\,\dd v .
\]
\end{theorem}
\admitted

\begin{remark}
A demonstração completa --- aproximar $\Phi$ por sua \omterm{def:b2:diffcalc:differential}{diferencial} numa
grade fina e controlar as células do bordo --- é longa, embora
não profunda; ela é feita por inteiro na teoria da medida do terceiro ano, como
consequência da teoria de Lebesgue. A heurística é a imagem
já usada para a \omterm{def:b2:multint:domain}{área} de superfícies: um pequeno quadrado de lado $\dd u$ em
$(u, v)$ é aplicado, em primeira ordem, sobre o paralelogramo \omterm{def:b2:structures:generated}{gerado}
por $\Phi_u\,\dd u$ e $\Phi_v\,\dd v$, cuja \omterm{def:b2:multint:domain}{área} é
$\abs{\det J_\Phi}\,\dd u\,\dd v$
(\cref{lem:b2:surfaces:lagrange}).
\end{remark}

\begin{remark}[Método: escolhendo a mudança de variáveis]
Três reflexos cobrem a maioria dos casos. \emph{Simetria do
integrando}: $x^2 + y^2$ pede polares, uma estrutura de produto
pede que se mantenham os eixos cartesianos. \emph{Forma do
bordo}: bordos $u(x,y) = c_1$, $v(x,y) = c_2$ imploram pelas
próprias coordenadas $(u, v)$, como no exemplo da região
hiperbólica abaixo --- o domínio se torna um retângulo,
o que já é toda a vitória. \emph{Estrutura linear}:
expressões em $x + y$ e $x - y$ convidam à rotação de
$45$ graus ou a um cisalhamento (\cref{ex:b2:multint:affine}). Em todos
os casos, três caixas a marcar antes de integrar: a aplicação é uma
bijeção do novo domínio sobre o antigo; seu jacobiano
é calculado \emph{no sentido efetivamente usado} (inverta no
fim, se for mais fácil); e o jacobiano entra com seu
valor absoluto.
\end{remark}

\begin{example}[Coordenadas polares]\label{ex:b2:multint:polar}
$\Phi(\rho, \alpha) = (\rho\cos\alpha,\ \rho\sin\alpha)$ tem
\[
J_\Phi = \begin{pmatrix}
\cos\alpha & -\rho\sin\alpha\\
\sin\alpha & \rho\cos\alpha
\end{pmatrix},
\qquad \det J_\Phi = \rho ,
\]
logo $\dd x\,\dd y = \rho\,\dd\rho\,\dd\alpha$. Para o disco $D_R$
de raio $R$:
\[
\iint_{D_R} e^{-(x^2 + y^2)}\,\dd x\,\dd y
= \int_0^{2\pi}\!\!\int_0^R e^{-\rho^2}\rho\,\dd\rho\,\dd\alpha
= \pi\bigl(1 - e^{-R^2}\bigr)
\xrightarrow[R\to\infty]{} \pi .
\]
Comparando com o quadrado $[-R, R]^2$ (que fica espremido entre
os discos $D_R$ e $D_{R\sqrt2}$, todos os integrandos positivos) obtém-se
$\bigl(\int_{-\infty}^\infty e^{-x^2}\dd x\bigr)^2 = \pi$:
\[
\boxed{\ \int_{-\infty}^{+\infty} e^{-x^2}\,\dd x = \sqrt{\pi}\ }
\]
--- de novo a integral de Gauss, agora por sua demonstração mais famosa
(compare com a dedução a uma variável em
\cref{ch:b2:integration}).
\end{example}

\begin{example}[Mudanças de variáveis afins]\label{ex:b2:multint:affine}
Para uma \omterm{def:b2:affine:subspace}{aplicação afim} $\Phi(u, v) = M(u, v)^{\mathsf T} + C$ com
$M$ invertível, o jacobiano é a matriz constante $M$:
as \omterm{def:b2:multint:domain}{áreas} são multiplicadas pelo fator constante $\abs{\det M}$
--- a promessa feita no~\cref{ch:b2:affine} é agora um teorema.
Dois usos imediatos. A elipse $\frac{x^2}{a^2} +
\frac{y^2}{b^2} \leq 1$ é a imagem do disco unitário por
$(u, v) \mapsto (au, bv)$, de modo que sua \omterm{def:b2:multint:domain}{área} é $ab \cdot \pi$ ---
sem cálculo algum. E, para a integral de $f(x + y)$ sobre o
quadrado $K = \intcc01^2$, o cisalhamento $\Phi(u, v) = (u - v, v)$
(\omterm{def:b2:linalg:det}{determinante} $1$) a transforma numa integral de $f(u)$ sobre um
paralelogramo, que Fubini fatia a $u$ constante: com $f =
\exp$,
\[
\iint_K \eu^{x+y}\,\dd x\,\dd y
= \Bigl(\int_0^1 \eu^x\,\dd x\Bigr)^2 = (\eu - 1)^2,
\]
como a estrutura de produto confirma. Escolher coordenadas
adaptadas ao integrando --- e não ao domínio --- é a
outra metade do ofício.
\end{example}


\begin{example}[Coordenadas adaptadas a um domínio
curvilíneo]\label{ex:b2:multint:hyperbolicregion}
Seja $D$ a região do primeiro quadrante limitada pelas
hipérboles $xy = 1$ e $xy = 3$ e pelas retas $y = x$ e
$y = 3x$. Nas coordenadas $u = xy$, $v = y/x$ o domínio
se torna o quadrado $\intcc13 \times \intcc13$; invertendo,
\[
x = \sqrt{u/v}, \qquad y = \sqrt{uv},
\qquad
\det J = x_uy_v - x_vy_u = \frac{1}{2v}
\]
(um cálculo de duas linhas com $x = u^{1/2}v^{-1/2}$, $y =
u^{1/2}v^{1/2}$). Portanto
\[
\operatorname{Area}(D)
= \int_1^3\!\!\int_1^3\frac{\dd u\,\dd v}{2v}
= 2\cdot\frac{\ln 3}{2} = \ln 3 \approx 1.10 .
\]
Tentar fatiar $D$ em coordenadas cartesianas significa recortá-lo
em três pedaços com bordos hiperbólicos e lineares ---
factível, sem alegria e propenso a erros. A moral repete a
\cref{ex:b2:multint:affine} com força total: leia as
equações do bordo e deixe que \emph{elas} escolham as
coordenadas; o jacobiano então converte a \omterm{def:b2:multint:domain}{área} da célula da
malha curvilínea, exatamente como o $\rho$ fez para as coordenadas
polares.
\end{example}

\begin{example}[Valores médios]\label{ex:b2:multint:average}
O \emph{valor médio} de $f$ num domínio $D$ é
$\frac1{\operatorname{Area}(D)}\iint_Df$. Amostra: a distância
média ao centro para um ponto escolhido uniformemente no
disco de raio $R$ é
\[
\frac{1}{\pi R^2}\int_0^{2\pi}\!\!\int_0^R
\rho\cdot\rho\,\dd\rho\,\dd\alpha
= \frac{2\pi R^3/3}{\pi R^2} = \frac{2R}3 ,
\]
e não $R/2$: a \omterm{def:b2:multint:domain}{área} uniforme põe mais massa nos raios grandes (a
coroa de raio $\rho$ tem peso proporcional a $\rho$),
de modo que a média fica além da metade. Acertar esse
fator é exatamente o jacobiano polar em ação, e a
mesma ponderação explica o centroide $\bar z = 3R/8$ da
meia-bola calculado adiante no capítulo, em vez de $R/2$.
\end{example}

\section{O teorema de Green--Riemann}

\begin{theorem}[Green--Riemann]\label{thm:b2:multint:green}
Seja $K \subseteq \R^2$ um domínio \omterm{def:b2:metric:compact}{compacto} elementar nos
dois sentidos (ou uma união finita de tais domínios colados ao longo de segmentos),
com bordo $\partial K$ uma curva fechada $\mathcal{C}^1$ por pedaços
orientada no sentido \emph{anti-horário} (o domínio fica à
esquerda). Para $P, Q$ de classe $\mathcal{C}^1$ numa vizinhança de
$K$:
\[
\oint_{\partial K} P\,\dd x + Q\,\dd y
= \iint_K \Bigl(\frac{\partial Q}{\partial x}
- \frac{\partial P}{\partial y}\Bigr)\,\dd x\,\dd y .
\]
\end{theorem}

\begin{proof}
Primeiro, ambos os membros são aditivos quando se corta $K$ ao longo de um
segmento em dois pedaços $K_1, K_2$: as integrais duplas se somam
pela aditividade de $\iint$; quanto às integrais de bordo, os
bordos anti-horários de $K_1$ e de $K_2$ percorrem cada um
o corte interior uma vez, em sentidos \emph{opostos}, de modo que, na
soma
\[
\oint_{\partial K_1} + \oint_{\partial K_2}
= \oint_{\partial K} + (\text{corte, nos dois sentidos})
= \oint_{\partial K},
\]
as duas passagens ao longo do corte se cancelam
(\cref{prop:b2:multint:lineinv}) e só o bordo exterior
sobrevive. Iterando finitos cortes, basta tratar
um domínio elementar. Demonstramos $\oint P\,\dd x = -\iint_K P_y$ num
domínio $D = \{a \leq x \leq b,\ \varphi_1(x) \leq y
\leq \varphi_2(x)\}$ elementar em $y$; a identidade $\oint Q\,\dd y = \iint_K Q_x$
é \omterm{def:b2:quadratic:adjoint}{simétrica} (elementar em $x$), e o teorema é a soma das duas.

Calcule a integral dupla por Fubini e pelo teorema fundamental
a uma variável:
\[
\iint_D \frac{\partial P}{\partial y}\,\dd x\,\dd y
= \int_a^b \bigl(P(x, \varphi_2(x)) - P(x, \varphi_1(x))\bigr)
\,\dd x .
\]
Ora, o bordo de $D$, no sentido anti-horário, consiste em: o gráfico
inferior $y = \varphi_1(x)$ percorrido da esquerda para a direita, o segmento
vertical direito $x = b$ (para cima), o gráfico superior $y =
\varphi_2(x)$ percorrido \emph{da direita para a esquerda}, o segmento
vertical esquerdo $x = a$ (para baixo). Ao longo dos segmentos verticais $x$ é
constante, de modo que eles contribuem com $0$ para $\oint P\,\dd x$; os gráficos,
parametrizados por $x$, dão
\[
\oint_{\partial D} P\,\dd x
= \int_a^b P(x, \varphi_1(x))\,\dd x
- \int_a^b P(x, \varphi_2(x))\,\dd x
= -\iint_D \frac{\partial P}{\partial y}\,\dd x\,\dd y .
\qedhere
\]
\end{proof}

\begin{corollary}[Área pelo bordo]\label{cor:b2:multint:area}
Sob as hipóteses do~\cref{thm:b2:multint:green},
\[
\operatorname{Area}(K)
= \oint_{\partial K} x\,\dd y
= -\oint_{\partial K} y\,\dd x
= \frac12\oint_{\partial K} x\,\dd y - y\,\dd x .
\]
\end{corollary}

\begin{proof}
Aplique Green--Riemann a $(P, Q) = (0, x)$, $(-y, 0)$ e
$\frac12(-y, x)$: de cada vez, $Q_x - P_y = 1$.
\end{proof}

\begin{remark}[Escolhendo entre as três fórmulas de área]
As três fórmulas de bordo são iguais, mas não intercambiáveis na
prática. Use $\oint x\,\dd y$ quando a parametrização torna
$\dd y$ simples (gráficos sobre o eixo $y$), $-\oint y\,\dd x$
simetricamente, e a semissoma \omterm{def:b2:quadratic:adjoint}{simétrica} quando a
parametrização trata $x$ e $y$ de modo equilibrado --- para a elipse
ela produziu um integrando constante, sem nenhuma
linearização trigonométrica. Em bordos poligonais a semissoma
se torna a fórmula do cadarço do~\cref{exo:b2:multint:12},
o algoritmo dos agrimensores. E quando o bordo é percorrido
no sentido horário pela parametrização dada, as três fórmulas
devolvem \emph{menos} a \omterm{def:b2:multint:domain}{área}: um resultado negativo não é um
erro de cálculo, mas um relatório de orientação --- inverta o
sinal, ou a parametrização.
\end{remark}

\begin{example}[Área da elipse]\label{ex:b2:multint:ellipsearea}
Para $x = a\cos t$, $y = b\sin t$, $t \in [0, 2\pi]$:
\[
\operatorname{Area}
= \frac12\int_0^{2\pi}\bigl(a\cos t \cdot b\cos t
- b\sin t\cdot(-a\sin t)\bigr)\,\dd t
= \frac{ab}{2}\int_0^{2\pi}\dd t = \pi ab .
\]
\end{example}

\begin{example}[Green--Riemann como
verificação cruzada]\label{ex:b2:multint:greencheck}
Tome $P = -y^3$, $Q = x^3$ no disco unitário fechado $D$.
Lado do bordo, com $\gamma(t) = (\cos t, \sin t)$:
\[
\oint_{\partial D}P\,\dd x + Q\,\dd y
= \int_0^{2\pi}\bigl(\sin^4 t + \cos^4 t\bigr)\dd t
= 2\pi\cdot\Bigl(\frac38 + \frac38\Bigr) = \frac{3\pi}2 ,
\]
por linearização ($\sin^4 + \cos^4 = \tfrac34 +
\tfrac14\cos4t$). Lado do interior:
\[
\iint_D(Q_x - P_y)\,\dd x\,\dd y
= \iint_D 3(x^2 + y^2)\,\dd x\,\dd y
= 3\int_0^{2\pi}\!\!\int_0^1\rho^3\,\dd\rho\,\dd\alpha
= \frac{3\pi}2 .
\]
Mesmo número, dois cálculos muito diferentes --- e esse é
o uso prático: o lado da identidade de Green que for mais
fácil torna-se o cálculo, o outro uma verificação. Para
circulações de campos polinomiais em torno de curvas fechadas, a
integral dupla é quase sempre o lado fácil.
\end{example}

\begin{remark}
Green--Riemann explica o~\cref{ex:b2:multint:angleform}: para uma
\omterm{def:b2:multint:exact}{forma fechada} ($Q_x = P_y$), a integral em torno do bordo de qualquer
domínio contido em $U$ se anula. A forma ângulo deixa de ser \omterm{def:b2:multint:exact}{exata}
apenas porque o furo na origem impede que o disco limitado
pelo círculo unitário esteja dentro de $U$ --- as \omterm{def:b2:multint:lineint}{integrais de linha} de
\omterm{def:b2:multint:exact}{formas fechadas} detectam os buracos do domínio. (Levada mais longe,
essa observação se torna a cohomologia de de Rham.)
\end{remark}

\section{Integrais triplas}

A teoria se estende a três variáveis sem ideia nova alguma: Fubini
reduz $\iiint$ a três integrais a uma variável (seja por
\emph{fatiamento}: $\iiint_K f = \int\bigl(\iint_{K_z}
f\bigr)\dd z$ sobre as fatias horizontais $K_z$, seja por
\emph{empilhamento}: integrando primeiro em $z$ ao longo de varetas verticais),
e a fórmula de mudança de variáveis vale com o jacobiano
$3 \times 3$.

\begin{example}[Coordenadas cilíndricas e esféricas]\label{ex:b2:multint:spherical}
\emph{Cilíndricas} $(x, y, z) = (\rho\cos\alpha, \rho\sin\alpha,
z)$: $\dd x\,\dd y\,\dd z = \rho\,\dd\rho\,\dd\alpha\,\dd z$.
\emph{Esféricas} $(x, y, z) = (r\cos\theta\cos\varphi,\
r\sin\theta\cos\varphi,\ r\sin\varphi)$ ($\theta$ a longitude,
$\varphi \in [-\frac\pi2, \frac\pi2]$ a latitude): desenvolvendo o \omterm{def:b2:linalg:det}{determinante} $3
\times 3$ pela última linha,
\[
\det J = r^2\cos\varphi ,
\qquad
\dd x\,\dd y\,\dd z
= r^2\cos\varphi\;\dd r\,\dd\theta\,\dd\varphi .
\]
\omterm{pb:b2:multint:1}{Volume da bola} de raio $R$:
\[
V = \int_0^R\!\!\int_0^{2\pi}\!\!\int_{-\pi/2}^{\pi/2}
r^2\cos\varphi\;\dd\varphi\,\dd\theta\,\dd r
= \frac{R^3}{3}\cdot 2\pi \cdot 2
= \boxed{\frac43\pi R^3} ,
\]
quitando enfim a fórmula admitida nos capítulos de volume dos
livros anteriores.
\end{example}

\begin{example}[O tetraedro, duas vezes]
\label{ex:b2:multint:tetra}
O volume de $T = \{x, y, z \geq 0,\ x + y + z \leq 1\}$, por
empilhamento: para $(x, y)$ fixo no triângulo $x + y \leq 1$,
$z$ percorre $\intcc0{1 - x - y}$, logo
\[
V = \int_0^1\!\!\int_0^{1-x}(1 - x - y)\,\dd y\,\dd x
= \int_0^1\frac{(1 - x)^2}{2}\,\dd x = \frac16 .
\]
Por fatiamento: a seção à altura $z$ é o triângulo $\{x, y
\geq 0,\ x + y \leq 1 - z\}$, de \omterm{def:b2:multint:domain}{área} $\frac{(1-z)^2}2$, e
$V = \int_0^1\frac{(1-z)^2}2\,\dd z = \frac16$ de novo --- os
dois cálculos são as mesmas integrais em ordem
diferente, que é tudo o que Fubini afirma. O valor $\frac16 =
\frac13\cdot\frac12\cdot1$ é a fórmula do cone
(\cref{ex:b2:multint:cone}) com base triangular, e a versão
em dimensão $n$, $1/n!$, é demonstrada exatamente por esse
fatiamento no problema de fim de semana.
\end{example}

\begin{example}[Centroide de uma meia-bola]\label{ex:b2:multint:centroid}
Para a meia-bola superior $H$ de raio $R$ ($z \geq 0$), a
altura do centroide é $\bar z = \frac1{V}\iiint_H z$, com $V =
\frac23\pi R^3$. Em coordenadas esféricas ($z =
r\sin\varphi$, $\varphi \in \intcc0{\pi/2}$):
\[
\iiint_H z
= \int_0^R r^3\,\dd r\int_0^{2\pi}\dd\theta
\int_0^{\pi/2}\sin\varphi\cos\varphi\,\dd\varphi
= \frac{R^4}4\cdot2\pi\cdot\frac12 = \frac{\pi R^4}4 ,
\]
logo
\[
\bar z = \frac{\pi R^4/4}{2\pi R^3/3} = \frac{3R}8 :
\]
o ponto de equilíbrio de um hemisfério maciço fica a três oitavos
do raio acima da face plana --- abaixo da meia-altura
$R/2$, como deve ser, pois o sólido é mais gordo perto da base.
Todo cálculo de centroide tem essa forma: uma integral de
momento, um volume, uma razão e uma verificação de plausibilidade
contra a geometria.
\end{example}

\begin{example}[Volume por fatiamento: o cone]\label{ex:b2:multint:cone}
Um cone de \omterm{def:b2:multint:domain}{área} de base $A$ e altura $h$ (vértice em cima, base em $z = 0$):
a fatia à altura $z$ é a base escalada pelo fator $(1 -
z/h)$, de \omterm{def:b2:multint:domain}{área} $A(1 - z/h)^2$. Portanto
\[
V = \int_0^h A\Bigl(1 - \frac zh\Bigr)^2\dd z = \frac{Ah}{3} :
\]
o um terço das fórmulas da escola, válido para \emph{qualquer} forma
de base --- o fatiamento o transforma na integral de um quadrado.
\end{example}

\begin{example}[Limiares de integrabilidade no
plano]\label{ex:b2:multint:threshold}
Para quais $\alpha > 0$ a integral $\iint_{D}\rho^{-\alpha}\,\dd
x\,\dd y$ converge no disco unitário perfurado $D$ (limite
sobre coroas $\varepsilon \leq \rho \leq 1$)? Em coordenadas
polares,
\[
\int_0^{2\pi}\!\!\int_\varepsilon^1\rho^{-\alpha}\,
\rho\,\dd\rho\,\dd\alpha
= 2\pi\int_\varepsilon^1\rho^{1-\alpha}\,\dd\rho ,
\]
que converge quando $\varepsilon \to 0$ se e somente se $1 - \alpha > -1$,
isto é, $\alpha < 2$: em dimensão $2$ o expoente crítico de
singularidade é a própria dimensão, o $\rho$ suplementar vindo do
jacobiano suavizando a singularidade em uma potência. (Do mesmo modo,
$\alpha < 3$ para uma singularidade \omterm{def:b2:funcseq:def}{pontual} no espaço, via $r^2$.)
Contabilidade radial desse tipo é como a integrabilidade é
decidida num relance no quadro de Lebesgue do terceiro ano ---
e é a razão pela qual $\iiint 1/r$ convergiu sem esforço no
\cref{exo:b2:multint:7}.
\end{example}

\begin{remark}[Armadilhas comuns]
(i) \emph{Orientação}: uma \omterm{def:b2:multint:lineint}{integral de linha} muda de sinal com o
sentido de percurso, e Green--Riemann exige o bordo
no sentido anti-horário (domínio à esquerda); para um domínio com um
buraco, o bordo interno é percorrido no sentido \emph{horário}. (ii)
\emph{O jacobiano entra com valor absoluto}: a mudança de
variáveis nunca produz \omterm{def:b2:multint:domain}{área} negativa, e esquecer
$\abs{\det}$ tipicamente inverte sinais exatamente quando a aplicação
inverte a orientação. (iii) \emph{O fator polar $\rho$}:
$\dd x\,\dd y = \rho\,\dd\rho\,\dd\alpha$, e não
$\dd\rho\,\dd\alpha$ --- o erro mais comum de todo o
capítulo; a análise dimensional o pega, pois
$\dd\rho\,\dd\alpha$ tem dimensão de comprimento, não de
\omterm{def:b2:multint:domain}{área}. (iv) \emph{Integrais duplas impróprias}: os limites sobre
discos crescentes e sobre quadrados crescentes coincidem aqui porque os
integrandos são positivos (aperto); para integrandos que mudam de
sinal o limite pode depender da exaustão, e nenhuma
afirmação é feita sem convergência absoluta. (v)
\emph{Domínios contra integrandos}: um integrando produto num
domínio que não é produto \emph{não} fatora a integral ---
a fatoração precisa dos dois, como no quadrado do
\cref{ex:b2:multint:polar}.
\end{remark}

\begin{remark}[Perspectivas dentro deste volume]
A integral de Gauss aqui calculada está discretamente em toda parte nos
capítulos de probabilidade: a constante $\sqrt\pi$ dentro da
fórmula de Stirling (\cref{thm:b2:comparison:stirling}) é
a integral deste capítulo e, através de Stirling, ela fixa a
assintótica em $1/\sqrt{\pi n}$ das probabilidades de retorno do
passeio aleatório no~\cref{ch:b2:proba}. As \omterm{pb:b2:multint:1}{integrais de Wallis} do
problema de fim de semana reaparecem lá também, movendo as mesmas
estimativas do binomial central. Na outra direção, os elementos de
\omterm{def:b2:multint:domain}{área} e de volume deste capítulo completam a geometria do
\cref{ch:b2:surfaces}, e a fórmula de Green recalcula as \omterm{def:b2:multint:domain}{áreas} das
\omterm{pb:b2:curves:1}{envoltórias} do~\cref{ch:b2:curves} (o \omterm{pb:b2:curves:1}{astroide}, no
\cref{exo:b2:multint:5}). Um capítulo, três serviços:
medida para a geometria, constantes para a probabilidade e a
disciplina da mudança de variáveis usada pelas duas.
\end{remark}

\section{Exercícios}

\begin{exercise}[$\star$]\label{exo:b2:multint:1}
Calcule $\int_\gamma y^2\,\dd x + x\,\dd y$ ao longo de: (a) o segmento
de $(0,0)$ a $(1,1)$; (b) o arco de parábola $y = x^2$ de
$(0,0)$ a $(1,1)$. A forma é \omterm{def:b2:multint:exact}{exata}?
\end{exercise}

\begin{exercise}[$\star$]\label{exo:b2:multint:2}
Mostre que $\omega = (2xy + y^3)\,\dd x + (x^2 + 3xy^2 + 1)\,\dd y$
é fechada em $\R^2$, determine um \omterm{def:b2:multint:exact}{potencial} e calcule
$\int_\gamma\omega$ ao longo de qualquer arco de $(0, 0)$ a $(1, 2)$.
\end{exercise}

\begin{exercise}[$\star$]\label{exo:b2:multint:3}
Calcule $\iint_D (x + y)\,\dd x\,\dd y$ em que $D$ é o domínio
limitado por $y = x^2$ e $y = x$ ($0 \leq x \leq 1$), nas duas
\omterm{def:b2:structures:generated}{ordens} de integração.
\end{exercise}

\begin{exercise}[$\star\star$]\label{exo:b2:multint:4}
Usando coordenadas polares, calcule $\iint_D \frac{\dd x\,\dd y}{(1 +
x^2 + y^2)^2}$ em todo o plano (como limite sobre discos) e
$\iint_{D'} xy\,\dd x\,\dd y$ no quarto de disco $D' = \{x, y
\geq 0,\ x^2 + y^2 \leq 1\}$.
\end{exercise}

\begin{exercise}[$\star\star$]\label{exo:b2:multint:5}
Calcule a \omterm{def:b2:multint:domain}{área} encerrada pelo \omterm{pb:b2:curves:1}{astroide} $x = \cos^3 t$, $y =
\sin^3 t$, $t \in [0, 2\pi]$, usando o
\cref{cor:b2:multint:area}. \emph{(Linearize $\sin^2 t\cos^2 t$.)}
\end{exercise}

\begin{exercise}[$\star\star$]\label{exo:b2:multint:6}
Calcule o volume do sólido limitado abaixo pelo paraboloide $z
= x^2 + y^2$ e acima pelo plano $z = 1$, pelos dois métodos:
empilhamento (integre $1 - x^2 - y^2$ sobre o disco unitário, coordenadas
polares) e fatiamento (as fatias horizontais são discos de raio
$\sqrt z$).
\end{exercise}

\begin{exercise}[$\star\star$]\label{exo:b2:multint:7}
(Atração gravitacional de uma bola --- teorema de Newton, caso
particular) Mostre que o volume da casca esférica $a \leq r \leq
b$ é $\frac43\pi(b^3 - a^3)$ e calcule $\iiint_{B}
\frac{\dd x\,\dd y\,\dd z}{r}$ na bola $B$ de raio $R$
($r$ a distância à origem). \emph{(Coordenadas
esféricas.)}
\end{exercise}

\begin{exercise}[$\star\star\star$]\label{exo:b2:multint:8}
(Lema de Poincaré, caso estrelado) Seja $U$ estrelado em
relação a $0$ (isto é, $M \in U \Rightarrow [0, M] \subseteq U$)
e seja $\omega = P\dd x + Q\dd y$ uma forma $\mathcal{C}^1$ fechada em
$U$. Defina
\[
f(x, y) = \int_0^1 \bigl(x\,P(tx, ty) + y\,Q(tx, ty)\bigr)\dd t .
\]
Usando a derivação sob o sinal de integral
(\cref{ch:b2:integration}) e $P_y = Q_x$, mostre que $f_x = P$
e $f_y = Q$: toda \omterm{def:b2:multint:exact}{forma fechada} num \omterm{def:b2:metric:topology}{aberto} estrelado é
\omterm{def:b2:multint:exact}{exata}.
\end{exercise}

\begin{exercise}[$\star\star\star$]\label{exo:b2:multint:9}
(Integral de Dirichlet por integração dupla) Justifique e explore
\[
\int_0^\infty\!\!\int_0^\infty e^{-xy}\sin x\;\dd y\,\dd x
\quad\text{vs}\quad
\int_0^\infty\!\!\int_0^\infty e^{-xy}\sin x\;\dd x\,\dd y
\]
em $[0, A] \times [0, \infty)$: mostre que $\int_0^A \frac{\sin
x}{x}\dd x = \frac\pi2 - \int_0^\infty
e^{-Ay}\frac{y\sin A + \cos A}{1 + y^2}\dd y$ e recupere
$\int_0^\infty \frac{\sin x}{x}\,\dd x = \frac\pi2$, comparando
com a demonstração por \omterm{thm:b2:integration:continuity}{integral com parâmetro} do
\cref{ch:b2:integration}.
\end{exercise}

\begin{exercise}[$\star\star\star$]\label{exo:b2:multint:10}
(Desigualdade isoperimétrica via Wirtinger) Seja $\gamma$ uma curva
fechada simples $\mathcal{C}^1$ de \omterm{def:b2:curves:length}{comprimento} $2\pi$, parametrizada por \omterm{def:b2:curves:length}{comprimento
de arco} em $[0, 2\pi]$, encerrando \omterm{def:b2:multint:domain}{área} $A$. Usando
o~\cref{cor:b2:multint:area}, Parseval e a desigualdade de Wirtinger
(exercícios do~\cref{ch:b2:fourier}), demonstre que $A \leq \pi$, com
igualdade para o círculo. \emph{(Normalize $\int_0^{2\pi} x(s)\dd s
= 0$; escreva $2A = \oint x\,\dd y - y\,\dd x$ e majore $2A \leq
\int (x^2 + y'^2)$ com cuidado via $2A = \int_0^{2\pi}(xy' - yx')\dd
s$ e $x^2 + y'^2 \geq 2xy'$.)}
\end{exercise}

\begin{exercise}[$\star\star$]\label{exo:b2:multint:11}
(Momentos da bola) Para a bola $B$ de raio $R$ em
$\R^3$, calcule $\iiint_B z^2\,\dd x\,\dd y\,\dd z$ em
coordenadas esféricas e deduza $\iiint_B (x^2 + y^2 +
z^2)\,\dd x\,\dd y\,\dd z$ por simetria. Confira o
último resultado contra o cálculo por cascas $\int_0^R r^2\cdot4\pi
r^2\,\dd r$.
\end{exercise}

\begin{exercise}[$\star\star$]\label{exo:b2:multint:12}
(Fórmula do cadarço) Seja $K$ um polígono de vértices $(x_1,
y_1), \dots, (x_m, y_m)$ em ordem anti-horária (índices
módulo $m$). Deduza do~\cref{cor:b2:multint:area} que
\[
\operatorname{Area}(K)
= \frac12\sum_{i=1}^m
\bigl(x_iy_{i+1} - x_{i+1}y_i\bigr) ,
\]
e confira a fórmula no triângulo $(0,0)$, $(1,0)$,
$(0,1)$.
\end{exercise}

\section{Problema: o volume da bola em dimensão $n$}

\begin{problem}[{Problema de fim de semana --- $V_n =
\pi^{n/2}/\Gamma(\frac n2 + 1)$, e a estranheza das
dimensões altas}]\label{pb:b2:multint:1}
O disco tem \omterm{def:b2:multint:domain}{área} $\pi$, o \omterm{pb:b2:multint:1}{volume da bola} vale $\frac43\pi$ --- e
depois? Este problema calcula o volume da bola unitária de
$\R^n$ para todo $n$, duas vezes (por uma recursão de fatiamento movida
pelas \omterm{pb:b2:multint:1}{integrais de Wallis}, depois via a função $\Gamma$ e a
integral de Gauss do~\cref{ex:b2:multint:polar}), e depois lê
a geometria: os volumes atingem o máximo na dimensão cinco e correm
para zero, e quase toda a bola de dimensão alta se esconde numa
casca fina perto de seu bordo.\index{volume da bola}\index{integrais
de Wallis} Para uma função \omterm{def:b2:metric:continuity}{contínua} numa bola de $\R^n$, a
integral é entendida como a integral iterada $n$ vezes
(fatiando uma coordenada de cada vez, como no capítulo para $n
\leq 3$); escrevemos $B_n(R)$ para a bola fechada de raio $R$
centrada em $0$, $v_n(R)$ para seu volume e $V_n =
v_n(1)$, com $V_0 = 1$ por convenção.

\medskip
\noindent\textbf{Parte I --- A recursão de fatiamento.}
\begin{enumerate}
  \item Substituindo $x_i = Ru_i$ em cada uma das $n$ integrais
        iteradas, mostre que $v_n(R) = V_nR^n$.
  \item Fatiando $B_n(1)$ ao longo de sua última coordenada, mostre que
        \[
        V_n = V_{n-1}\int_{-1}^{1}(1 -
        t^2)^{\frac{n-1}2}\,\dd t .
        \]
  \item Com $t = \sin\theta$, identifique a integral como uma
        \omterm{pb:b2:multint:1}{integral de Wallis}: $\int_{-1}^1(1 -
        t^2)^{\frac{n-1}2}\dd t = 2W_n$, em que $W_n =
        \int_0^{\pi/2}\cos^n\theta\,\dd\theta =
        \int_0^{\pi/2}\sin^n\theta\,\dd\theta$.
  \item Demonstre as duas identidades de Wallis (integre por partes;
        depois telescope $nW_nW_{n-1}$):
        \[
        W_n = \frac{n-1}nW_{n-2} \quad (n \geq 2),
        \qquad
        W_nW_{n-1} = \frac{\pi}{2n} \quad (n \geq 1).
        \]
\end{enumerate}

\medskip
\noindent\textbf{Parte II --- A recursão resolvida.}
\begin{enumerate}[resume]
  \item Combine as questões 2--4 na recursão de dois passos
        \[
        V_n = \frac{2\pi}{n}\,V_{n-2}
        \qquad (n \geq 2).
        \]
  \item Deduza as \omterm{def:b2:multint:exact}{formas fechadas}, para $k \geq 0$:
        \[
        V_{2k} = \frac{\pi^k}{k!},
        \qquad
        V_{2k+1} = \frac{2^{k+1}\pi^k}{1\cdot3\cdot5\cdots
        (2k+1)} .
        \]
  \item Tabule $V_1, \dots, V_7$ numericamente. Usando a
        razão $V_n/V_{n-2} = 2\pi/n$ e os valores de $2W_5$
        e de $2W_6$, demonstre que a sequência $(V_n)$ cresce
        até seu máximo $V_5 = \frac{8\pi^2}{15} \approx
        5.26$ e decresce a partir daí.
  \item Mostre que $V_n \to 0$ mais rápido que qualquer sequência
        geométrica, e que $\sum_{n\geq1} V_n$ converge: todas
        as bolas unitárias juntas têm volume total finito.
  \item Demonstre a identidade geradora
        \[
        \sum_{k\geq0} V_{2k}\,x^{2k} = \eu^{\pi x^2}
        \qquad (x \in \R),
        \]
        e deduza $\sum_{k\geq0}V_{2k} = \eu^\pi \approx
        23.14$.
\end{enumerate}

\medskip
\noindent\textbf{Parte III --- Segunda via: $\Gamma$ e a
integral de Gauss.}
\begin{enumerate}[resume]
  \item Mostre, por Fubini (o integrando é um produto), que
        \[
        I_n = \int_{\R^n}\eu^{-\norm
        x^2}\dd x
        = \Bigl(\int_{-\infty}^{+\infty}
        \eu^{-t^2}\dd t\Bigr)^{\!n} = \pi^{n/2},
        \]
        a integral de Gauss em dimensão $n$, entendida como
        limite sobre cubos $\intcc{-R}{R}^n$.
  \item Recorde $\Gamma(s) = \int_0^\infty t^{s-1}\eu^{-t}\dd
        t$ (\cref{def:b2:integration:gamma}). A partir de
        $\Gamma(s+1) = s\,\Gamma(s)$
        (\cref{thm:b2:integration:gammaprops}) e de
        $\Gamma(\tfrac12) = \sqrt\pi$ (substitua $t = u^2$
        e invoque a integral de Gauss), calcule
        \[
        \Gamma(k + 1) = k!,
        \qquad
        \Gamma\Bigl(k + \frac32\Bigr) =
        \frac{1\cdot3\cdots(2k+1)}{2^{k+1}}\,\sqrt\pi .
        \]
  \item Demonstre, por indução através da recursão da questão
        5, a fórmula única
        \[
        V_n = \frac{\pi^{n/2}}{\Gamma\bigl(\frac n2 +
        1\bigr)} \qquad (n \geq 1),
        \]
        e verifique que ela reproduz ambas as \omterm{def:b2:multint:exact}{formas fechadas} da questão
        6.
  \item Mostre que $\int_0^\infty \eu^{-r^2}r^{n-1}\dd r =
        \tfrac12\Gamma\bigl(\tfrac n2\bigr)$ e deduza a
        identidade
        \[
        I_n = n\,V_n\int_0^\infty \eu^{-r^2}\,r^{n-1}\,\dd r.
        \]
        Interprete-a: a massa gaussiana de $\R^n$ é
        recolhida ao longo de cascas esféricas cuja
        ``\omterm{def:b2:multint:domain}{área} em dimensão $(n-1)$'' no raio $r$ vale
        $nV_nr^{n-1}$ --- ambos os membros estão agora demonstrados
        independentemente, de modo que a interpretação não custa nada.
  \item Ponha $s_{n-1} = nV_n$ (a \omterm{def:b2:multint:domain}{área} da esfera unitária
        $S^{n-1}$, coerentemente com $v_n(R) = \int_0^R
        s_{n-1}r^{n-1}\dd r$). Tabule $s_0, \dots, s_3$ e
        verifique $s_1 = 2\pi$, $s_2 = 4\pi$, $s_3 = 2\pi^2$.
\end{enumerate}

\medskip
\noindent\textbf{Parte IV --- As dimensões altas são estranhas.}
\begin{enumerate}[resume]
  \item A partir da fórmula de Stirling
        (\cref{thm:b2:comparison:stirling}) aplicada a $k!$,
        mostre que, para $n = 2k$ par:
        \[
        V_n \sim \frac{1}{\sqrt{\pi n}}
        \Bigl(\frac{2\pi\eu}{n}\Bigr)^{n/2}
        \qquad (n \to \infty, \ n \text{ par}),
        \]
        e explique por que a mesma cota de decaimento supergeométrico
        se estende aos $n$ ímpares via a recursão.
  \item A bola unitária está no cubo $\intcc{-1}1^n$ de
        volume $2^n$. Calcule a razão de preenchimento $V_n/2^n$ para
        $n = 2, 3, 10$ e mostre que ela tende a $0$: em dimensão
        alta, essencialmente todo o cubo está em seus
        cantos.
  \item Mostre que a fração de $v_n(1)$ situada a
        distância no máximo $\varepsilon$ da esfera de bordo é $1 -
        (1 - \varepsilon)^n \to 1$; numericamente, que
        fração de uma bola de dimensão $100$ está na
        casca externa de espessura $1\%$?
  \item Demonstre a assintótica de Wallis $W_n \sim
        \sqrt{\dfrac{\pi}{2n}}$ \emph{(monotonia de
        $(W_n)$, a razão $W_n/W_{n-2} \to 1$ e $W_nW_{n-1}
        = \frac\pi{2n}$)}, e a cota inferior $W_n \geq
        \sqrt{\dfrac{\pi}{2(n+1)}}$ para todo $n$.
  \item (Concentração numa faixa) A fração da bola unitária
        com primeira coordenada além de $\delta$ vale
        $\int_\delta^1(1 - x^2)^{\frac{n-1}2}\dd x \,\big/\,
        (2W_n)$. Usando $1 - u \leq \eu^{-u}$ e a cota
        de cauda $\int_\delta^\infty \eu^{-a x^2}\dd x \leq
        \frac{\eu^{-a\delta^2}}{2a\delta}$, mostre que essa
        fração é no máximo
        \[
        \frac{\eu^{-(n-1)\delta^2/2}}{(n-1)\,\delta}
        \Big/ \sqrt{\frac{2\pi}{n+1}}
        \]
        e conclua: para $\delta = s/\sqrt{n-1}$, tudo salvo uma
        fração $O(\eu^{-s^2/2}/s)$ da bola está na
        faixa $\abs{x_1} \leq s/\sqrt{n-1}$. Uma bola de dimensão
        alta é, estatisticamente, uma panqueca fina em todas
        as direções ao mesmo tempo.
  \item Reúna as questões 16--19 num parágrafo: onde
        está o volume de $B_n(1)$ (perto da esfera de bordo,
        e no entanto dentro de faixas $O(1/\sqrt n)$ de todo
        hiperplano pelo centro), e por que essas duas
        afirmações não se contradizem.
\end{enumerate}

\medskip
\noindent\textbf{Parte V --- Outros corpos, e síntese.}
\begin{enumerate}[resume]
  \item (Simplexo) Seja $\Delta_n = \{x \in \R^n : x_i \geq 0,\
        \sum x_i \leq 1\}$. Mostre, por fatiamento e indução, que
        $\operatorname{vol}(\Delta_n) = \frac1{n!}$.
  \item (Politopo cruzado) Deduza que $C_n = \{x :
        \sum\abs{x_i} \leq 1\}$ tem volume $\frac{2^n}{n!}$
        e verifique o sanduíche $C_n \subseteq B_n(1)
        \subseteq \intcc{-1}1^n$ ao nível dos volumes:
        $\frac{2^n}{n!} \leq V_n \leq 2^n$.
  \item Calcule $V_4$ de uma terceira maneira: fatie $\R^4 = \R^2 \times
        \R^2$, integre a \omterm{def:b2:multint:domain}{área} do disco em $(z, w)$ sobre
        o disco em $(x, y)$ em coordenadas polares e recupere
        $V_4 = \frac{\pi^2}2$.
  \item (Monte Carlo em apuros) Sorteia-se um ponto \omterm{def:b2:funcseq:def}{uniformemente}
        no cubo $\intcc{-1}1^{20}$. Mostre que a
        probabilidade de ele cair na bola inscrita é
        $V_{20}/2^{20} \approx 2.5\cdot10^{-8}$, de modo que cerca de
        quarenta milhões de sorteios são necessários antes que o primeiro acerto
        seja esperado: estimar $V_n$ por amostragem por rejeição
        desmorona em dimensão alta (a maldição da
        dimensionalidade).
  \item Síntese. Duas deduções independentes se encontraram em $V_n =
        \pi^{n/2}/\Gamma(\frac n2 + 1)$: liste qual teorema
        deste capítulo cada uma usou (Fubini, mudança de
        variáveis, a integral de Gauss em polares) e quais
        insumos a uma variável (Wallis, $\Gamma$, Stirling).
        Onde o volume do terceiro ano de graduação refaz esse cálculo
        com a teoria de Lebesgue, e o que ela acrescenta?
\end{enumerate}
\end{problem}
